\chapter{Reescritura DPO de operaciones POSIX}
\label{cap:dpo}

Este es el capítulo más cargado del libro y el corazón teórico de sotFS.
Las operaciones POSIX (\opcreate, \opmkdir, \opunlink, \oprename,
\oplink, \oprmdir, \opwrite) se modelan acá como reglas de
\emph{double-pushout} (DPO) sobre el TypeGraph del cap.~\ref{cap:typegraph}.
La virtud del enfoque DPO no es estética: cada regla viene con
condiciones formales (\emph{gluing condition}) cuya verificación es
mecánica, y la preservación de los invariantes de
sec.~\ref{sec:invariantes} se sigue del enunciado de la regla, no de un
chequeo aislado.

\section{Repaso: pushout en \texorpdfstring{$\Set$}{Set} y en \texorpdfstring{$\Graph$}{Graph}}

\begin{defn}[Pushout en $\Set$]
\label{def:pushout-set}
Dadas tres funciones $f: A \to B$, $g: A \to C$, su \emph{pushout} es un
objeto $D$ junto con dos funciones $h: B \to D$, $k: C \to D$ tales que
$h \circ f = k \circ g$ y, para cualquier otro objeto $D'$ con flechas
$h', k'$ que cumpla la misma ecuación, existe una única flecha
$u: D \to D'$ que las factoriza.
\end{defn}

Concretamente en $\Set$, el pushout se construye como el cociente
$(B \sqcup C)/\sim$ donde $f(a) \sim g(a)$ para todo $a \in A$. La
intuición: $A$ es el ``pegamento'' que identifica qué elementos de $B$ y
$C$ son ``el mismo elemento'' en el resultado.

\begin{figure}[h]
\centering
\begin{tikzpicture}[node distance=15mm and 22mm]
  \node (A) {$A$};
  \node[right=of A] (B) {$B$};
  \node[below=of A] (C) {$C$};
  \node[below=of B] (D) {$D$};
  \draw[arr] (A) -- (B) node[edgelbl, midway] {$f$};
  \draw[arr] (A) -- (C) node[edgelbl, midway] {$g$};
  \draw[arr] (B) -- (D) node[edgelbl, midway] {$h$};
  \draw[arr] (C) -- (D) node[edgelbl, midway] {$k$};
\end{tikzpicture}
\caption{Diagrama de pushout. $D$ es el menor objeto que cierra el
cuadrado conmutativo.}
\label{fig:pushout-set}
\end{figure}

En la categoría $\Graph$ (grafos dirigidos y morfismos de grafos) la
construcción es análoga: el pushout existe siempre y se construye nodo
por nodo, arista por arista, identificando los pares que vienen del
``pegamento'' $A$. Este es el sustrato técnico que permite hablar de
\emph{reemplazar un fragmento por otro pegándolo a lo que queda}.

\section{Una regla DPO}

\begin{defn}[Regla DPO]
\label{def:dpo-rule}
Una \emph{regla DPO} es un span
$\rho: \dpo{L}{K}{R}$ en la categoría $\Graph$ (o $\TGraph$, en nuestro
caso). Intuitivamente:
\begin{itemize}
  \item $L$ es el patrón \emph{a remover} (\emph{left-hand side}).
  \item $K$ es la parte \emph{glue} que se conserva tras la
    transformación.
  \item $R$ es el patrón \emph{a insertar} (\emph{right-hand side}).
\end{itemize}
La inclusión $K \hookrightarrow L$ y $K \hookrightarrow R$ codifica qué
sobrevive y qué se reemplaza.
\end{defn}

Un \emph{paso DPO} sobre un grafo \emph{host} $G$ requiere:

\begin{enumerate}
  \item Un \emph{matching} $m: L \to G$, esto es, un morfismo de
    grafos que ubica $L$ dentro de $G$.
  \item La \emph{gluing condition}: el matching no debe ``dejar
    aristas colgando'' al borrar lo que sólo pertenece a $L$. Más
    formalmente, todo nodo de $L$ con aristas adyacentes en $G$ que no
    estén en $m(L)$ debe estar en $K$ (es decir, debe sobrevivir).
  \item Construir el pushout complement $D$ tal que $L \leftarrow K
    \rightarrow D$ y $L \rightarrow G$ formen un cuadrado pushout.
  \item Construir el pushout final $H$ entre $D$ y $R$ a lo largo de
    $K$.
\end{enumerate}

\begin{figure}[h]
\centering
\begin{tikzpicture}[node distance=12mm and 18mm]
  \node (L) {$L$};
  \node[right=of L] (K) {$K$};
  \node[right=of K] (R) {$R$};
  \node[below=of L] (G) {$G$};
  \node[below=of K] (D) {$D$};
  \node[below=of R] (H) {$H$};
  \draw[arr] (K) -- (L) node[edgelbl, midway] {};
  \draw[arr] (K) -- (R) node[edgelbl, midway] {};
  \draw[arr] (L) -- (G) node[edgelbl, midway] {$m$};
  \draw[arr] (K) -- (D) node[edgelbl, midway] {};
  \draw[arr] (R) -- (H) node[edgelbl, midway] {};
  \draw[arr] (D) -- (G) node[edgelbl, midway] {};
  \draw[arr] (D) -- (H) node[edgelbl, midway] {};
  \node at ($(L)!0.5!(D)$) {\scriptsize PO};
  \node at ($(R)!0.5!(D)$) {\scriptsize PO};
\end{tikzpicture}
\caption{Diagrama de un paso DPO: dos cuadrados pushout (PO) compartiendo
el lado central $K \to D$. El de la izquierda \emph{borra} lo que está
en $L \setminus K$; el de la derecha \emph{inserta} lo que está en $R
\setminus K$.}
\label{fig:dpo-square}
\end{figure}

\begin{observacion}[En sotFS las reglas son funciones Rust, no estructuras explícitas]
\label{obs:dpo-rust}
El runtime no implementa un motor DPO genérico que tome un span $(L, K,
R)$ y haga el pushout complement: implementa cada operación POSIX como
una función Rust en \citaCrate{sotfs-ops} que, en orden, (i) ubica el
matching (\texttt{resolve\_name}, \texttt{get\_dir},
\texttt{get\_inode}); (ii) chequea la gluing condition con
\emph{pre-condiciones} explícitas; (iii) muta el grafo para producir
$H$. La equivalencia con la formulación categórica es por
\emph{inspección}: cada función se diseña para que el efecto neto sobre
el grafo coincida con el pushout final. La gluing condition aparece
como un \texttt{return Err(\dots)} antes de mutar; los chequeos
post-mutación son los siete invariantes (sec.~\ref{sec:invariantes}).
\end{observacion}

\section{Producciones DPO de las operaciones POSIX}

Acá vienen las siete producciones que cubren la mayoría del API POSIX.
Cada una se presenta con una caja \texttt{(L, K, R)}, su gluing
condition formal, la firma exacta de la función Rust en
\citaCrate{sotfs-ops} y el lemma de preservación. Las pruebas formales
están en Coq y se discuten en cap.~\ref{cap:coq}.

\subsection{\opcreate}

\paragraph{Intuición.} Crear un archivo regular bajo $d/n$. Aparece un
nuevo \NInode{} con \texttt{link\_count=1} y una arista \econtains{} con
nombre $n$ desde $d$ hasta el inode nuevo.

\paragraph{Producción $\rho_{\mathrm{create}}$.}

\begin{itemize}
  \item $L$: el directorio $d$ (un único nodo \NDir).
  \item $K = L$: el directorio sobrevive entero (no se borra nada).
  \item $R$: $d$ \textbf{más} un nuevo \NInode{} $i$ con
    \texttt{link\_count=1} y una arista \econtains{} de $d$ a $i$ con
    nombre $n$.
\end{itemize}

\paragraph{Gluing condition (NAC, \emph{negative application condition}).}
\[
  \neg\,\exists\,e \in E\,:\;
    \tau_E(e) = \econtains \;\wedge\; \mathrm{src}(e) = d \;\wedge\; e.\texttt{name} = n.
\]
``No existe ya una entrada con ese nombre en $d$.'' En POSIX corresponde
a \eexist.

\paragraph{Firma Rust} (\citaLine{sotfs-ops/src/lib.rs}{33}).

\begin{lstlisting}[language=rust]
pub fn create_file(
    g: &mut TypeGraph,
    parent_dir: DirId,
    name: &str,
    uid: u32, gid: u32,
    permissions: Permissions,
) -> Result<InodeId, GraphError>
\end{lstlisting}

\begin{lema}[Preservación bajo \opcreate]
\label{lem:create-preserva}
Si $G$ es well-formed y $\rho_{\mathrm{create}}$ se aplica con éxito
sobre $G$ produciendo $H$, entonces $H$ es well-formed.
\end{lema}

\begin{proof}[Esquema]
Los siete invariantes:
\begin{itemize}
  \item Link-count: el inode nuevo tiene \texttt{link\_count = 1} y
    exactamente una arista \econtains{} entrante con $\texttt{name} \neq
    \texttt{".."}$. Coincide.
  \item Unique-names: la NAC garantiza que $n$ no existía en $d$;
    seguimos sin duplicados.
  \item Dir-self-ref: no se tocan dirs distintos de $d$, y $d$ ya
    tenía su \texttt{.}.
  \item No-dangling: la arista nueva conecta $d$ (vivo) e $i$ (recién
    insertado, vivo).
  \item Block-refcount: no se tocan bloques.
  \item No-dir-cycles: no se inserta un \NDir, así que la sub-relación
    \econtains-restringida-a-dirs no cambia.
  \item Cap-monotonicity: no se tocan caps.
\end{itemize}
La prueba formal mecanizada vive en \citaCoq{DpoCreate.v}, theorem
\texttt{create\_preserves\_WellFormed} en línea 335 (cap.~\ref{cap:coq}).
\end{proof}

\subsection{\opmkdir}

\paragraph{Producción.}

\begin{itemize}
  \item $L$: el directorio padre $d_{\text{p}}$ (su \NDir{} y su \NInode).
  \item $K = L$.
  \item $R$: $d_{\text{p}}$ \textbf{más} un nuevo \NInode{} $i$
    (\texttt{vtype = Directory}, \texttt{link\_count=2}), un nuevo
    \NDir{} $d$ con \texttt{inode\_id = i}, y \emph{tres} aristas
    \econtains: $d_{\text{p}} \xrightarrow{n} i$ (la entrada nueva en
    el padre), $d \xrightarrow{\,.\,} i$ (la self-ref del nuevo dir),
    $d \xrightarrow{\,..\,} d_{\text{p}}.\texttt{inode\_id}$ (el back-pointer al padre).
\end{itemize}

\paragraph{Gluing condition.} Análoga a \opcreate: $\neg\,\exists\,e \in
E_d^{\econtains}$ con $e.\texttt{name} = n$.

\paragraph{Por qué \texttt{link\_count=2}.} Un directorio nuevo recibe
dos aristas \econtains{} entrantes: la del padre con nombre $n$, y la
\texttt{.} desde sí mismo. La invariante~\ref{inv:link-count} excluye
``..'' pero \emph{no} ``.''. El \texttt{link\_count} inicial es por
construcción 2.

\paragraph{Firma Rust} (\citaLine{sotfs-ops/src/lib.rs}{94}).

\begin{lstlisting}[language=rust]
pub fn mkdir(
    g: &mut TypeGraph,
    parent_dir: DirId,
    name: &str,
    uid: u32, gid: u32,
    permissions: Permissions,
) -> Result<CreateResult, GraphError>
\end{lstlisting}

Resultado: el \texttt{InodeId} y el \texttt{DirId} del directorio nuevo.

\subsection{\oplink} (hard link)

\paragraph{Producción.}

\begin{itemize}
  \item $L$: el directorio $d$ y el inode existente $i$.
  \item $K = L$.
  \item $R$: $L$ \textbf{más} una arista \econtains{} de $d$ a $i$ con
    nombre $n$. \emph{Adicionalmente}: el atributo $i.\texttt{link\_count}$ se
    incrementa en uno.
\end{itemize}

La actualización de \texttt{link\_count} es un \emph{cambio de
atributo}, no un cambio topológico; el formalismo DPO
\emph{atribuido} (sobre $\TGraph$, no $\Graph$) lo cubre con un morfismo
sobre los atributos.

\paragraph{Gluing condition.}
\[
  \neg\,\exists\,e \in E_d^{\econtains} : e.\texttt{name} = n
  \;\;\wedge\;\;
  i.\texttt{link\_count} < \texttt{LINK\_MAX}.
\]

La segunda cláusula evita overflow del contador y, más sutilmente,
acota el \emph{treewidth} del grafo
(\trampref{tr:link-treewidth} más abajo y cap.~\ref{cap:topologia}).
\texttt{LINK\_MAX = 65535} en
\citaLine{sotfs-graph/src/graph.rs}{28}.

\paragraph{Firma Rust} (\citaLine{sotfs-ops/src/lib.rs}{272}).

\subsection{\opunlink}

\paragraph{Producción.}

\begin{itemize}
  \item $L$: el directorio $d$, el inode $i$ y la arista \econtains{}
    $d \xrightarrow{n} i$.
  \item $K$: depende de $i.\texttt{link\_count}$.
    \begin{itemize}
      \item Si $i.\texttt{link\_count} = 1$, $K = \{d\}$ (el inode se
        irá).
      \item Si $i.\texttt{link\_count} > 1$, $K = \{d, i\}$ (el inode
        sobrevive con \texttt{link\_count} decrementado).
    \end{itemize}
  \item $R = K$ \textbf{más} actualización del atributo
    $i.\texttt{link\_count} \leftarrow i.\texttt{link\_count} - 1$ si $i \in K$.
\end{itemize}

\paragraph{Gluing condition.} Que la arista exista
($\exists\,e: \mathrm{src}(e)=d \wedge e.\texttt{name}=n$) y, si
$i.\texttt{vtype} = \texttt{Directory}$, que se invoque \oprmdir{} en
vez de \opunlink{} (POSIX prohíbe \opunlink{} sobre directorios; en
sotFS la verificación está en la firma).

\paragraph{Firma Rust} (\citaLine{sotfs-ops/src/lib.rs}{322}).

\begin{trampa}[Orphan blocks bajo \opunlink{} de hard-linked file]
\label{tr:orphan-blocks}
Cuando $i.\texttt{link\_count} = 1 \to 0$, el inode desaparece y sus
aristas \epointsto{} también. Pero los \NBlk{} que apuntaba podrían
quedar con \texttt{refcount} desactualizado si la regla no decrementa
$b.\texttt{refcount}$ por cada arista borrada. Versión vulnerable: la
regla decrementaba \texttt{refcount} sólo a la primera vista del bloque,
saltándose colisiones (mismo bloque apuntado dos veces desde el mismo
inode con distintos offsets — caso de archivo con bloques repetidos
internamente). Resultado: \invref{inv:block-refcount} violada.

\textbf{Fix.} Decrementar una vez por arista \epointsto, no una vez por
bloque distinto. \emph{Lección}: contar aristas, no objetos.
\end{trampa}

\subsection{\oprmdir}

\paragraph{Producción.}

\begin{itemize}
  \item $L$: el directorio padre $d_{\text{p}}$, el directorio $d$, el
    inode $i$ pareado a $d$, y \emph{tres aristas}: $d_{\text{p}}
    \xrightarrow{n} i$, $d \xrightarrow{\,.\,} i$, $d \xrightarrow{\,..\,}
    d_{\text{p}}.\texttt{inode\_id}$.
  \item $K = \{d_{\text{p}}\}$: sólo sobrevive el padre.
  \item $R = K$.
\end{itemize}

\paragraph{Gluing condition.} ``$d$ está vacío'': no hay aristas
\econtains{} salientes de $d$ excepto las dos auto-aristas (\texttt{.}
y \texttt{..}). En POSIX corresponde a \enotempty.

\paragraph{Firma Rust} (\citaLine{sotfs-ops/src/lib.rs}{198}).

\begin{observacion}[Por qué \oprmdir{} es la prueba más difícil en Coq]
\label{obs:rmdir-coq}
La preservación del \invref{inv:no-dir-cycles} bajo \oprmdir{} requiere
un argumento de invariante auxiliar (que \emph{no hay hard links a
directorios}). Esa propiedad no está formalizada en el modelo Records
de \citaCoq{SotfsGraph.v}, lo que deja cuatro \texttt{Admitted} en
\citaCoq{DpoRmdir.v}. El cap.~\ref{cap:coq} los lista uno por uno; el
ej.~\ref{ej:coq-extension} pide formalizarlos.
\end{observacion}

\subsection{\oprename}

Es la operación más rica del API POSIX: cubre tres casos según los
directorios involucrados sean iguales (\emph{same-dir rename}) o
distintos (\emph{cross-dir rename}), y según haya o no un destino
ya existente que reemplazar (la semántica POSIX exige que un destino
existente sea reemplazado atómicamente).

\paragraph{Same-dir, sin destino preexistente.}

\begin{itemize}
  \item $L$: el directorio $d$, el inode $i$, la arista $d \xrightarrow{n} i$.
  \item $K = \{d, i\}$.
  \item $R = K$ \textbf{más} la arista $d \xrightarrow{n'} i$.
\end{itemize}

Implementación: \texttt{rename\_same\_dir}
(\citaLine{sotfs-ops/src/lib.rs}{441}); en este caso particular se
reusa la misma \texttt{EdgeId} y se reescribe sólo el atributo
\texttt{name}, evitando alocar y liberar.

\paragraph{Cross-dir.} La regla es esencialmente
\opcreate(d',n') ; \opunlink(d,n) atómica, con dos chequeos extra:
\begin{itemize}
  \item Si $i$ es un directorio, \emph{no debe ser ancestro de $d'$}
    (evita ciclos en la jerarquía). Es una condición no local: requiere
    seguir las aristas \texttt{..} desde $d'$ hacia arriba y verificar
    que no aparece $i$.
  \item Si $i$ es un directorio, su \texttt{..} en $d$ se reescribe para
    apuntar al inode pareado de $d'$.
\end{itemize}

\paragraph{Reemplazo de destino existente.} POSIX exige que si $d'/n'$ ya
existe, \oprename{} \emph{lo reemplace atómicamente}. La regla DPO
correspondiente expone una arista \esupersedes{} desde el inode nuevo
hacia el viejo durante el paso intermedio, lo que permite que la
verificación TLA\textsuperscript{+} \citaTLA{sotfs\_transactions.tla}
detecte la atomicidad como single-step.

\paragraph{Firma Rust} (\citaLine{sotfs-ops/src/lib.rs}{404}).

\subsection{\opwrite}

\paragraph{Intuición.} Escribir bytes a un archivo. Concretamente, alocar
(o reusar) bloques que cubran el rango \texttt{[offset, offset+len)}
y mantener las aristas \epointsto{} con su \texttt{offset}.

\paragraph{Producción (versión simple, un único bloque).}

\begin{itemize}
  \item $L$: el inode $i$.
  \item $K = L$.
  \item $R$: $L$ \textbf{más} un nuevo \NBlk{} $b$ con
    \texttt{refcount=1} y una arista $i \xrightarrow{\epointsto, \text{off}=k} b$.
\end{itemize}

Si una arista \epointsto{} ya cubre el rango, se reusa o se splittea
según el tamaño. La versión completa se descompone en
\texttt{write\_block} (\citaLine{sotfs-ops/src/lib.rs}{626}) y
\texttt{write\_data} (\citaLine{sotfs-ops/src/lib.rs}{670}).

\section{Gluing condition vs. invariantes}

Es útil distinguir dos clases de chequeo:

\begin{enumerate}
  \item \textbf{Gluing condition (pre-mutación).} Se verifica
    \emph{antes} de aplicar la regla. Su falla resulta en un
    \texttt{Err(\dots)} que aborta sin tocar el grafo. Codifica el
    contrato \emph{POSIX-like} (\eexist, \enoent, \enotempty).
  \item \textbf{Invariante (post-mutación, eventual).} Se verifica
    \emph{después} de la mutación, en modo paranoid o test, mediante
    \texttt{check\_invariants()}. Su falla indica un \emph{bug en la
    regla}, no un input inválido.
\end{enumerate}

La idea de DPO es \emph{minimizar la segunda lista}: si las gluing
conditions están bien diseñadas, las reglas preservan los invariantes
\emph{por construcción}. Los Admitteds en Coq de cap.~\ref{cap:coq}
existen porque algunas gluing conditions no son aún tan estrictas como
podrían (concretamente: falta el invariante auxiliar
\texttt{NoHardLinkToDir}).

\section{Algoritmo genérico de aplicación de una regla DPO}

Aunque el runtime no implementa un motor genérico, la formulación
abstracta es útil para entender los Coq files de cap.~\ref{cap:coq}.

\begin{algorithm}[H]
\SetAlgoLined
\KwIn{Grafo host $G$, regla $\rho: \dpo{L}{K}{R}$, matching candidato $m: L \to G$}
\KwOut{Grafo $H$ resultado, o \textsc{Fail}}
\caption{\textsc{ApplyDPO}}
\label{alg:apply-dpo}

\If{$\neg\;\textsc{IsValidMatching}(m, L, G)$}{
  \Return{\textsc{Fail}}\;
}
\If{$\neg\;\textsc{CheckGluing}(L, K, m, G)$}{
  \Return{\textsc{Fail}}\;
}

$D \gets G \setminus m(L \setminus K)$\;

$H \gets D \cup R$\;
\If{$\neg\;\textsc{CheckInvariants}(H)$}{
  \Return{\textsc{Fail}}\;
}
\Return{$H$}\;
\end{algorithm}

La línea 4 (``borra lo que vive sólo en $L$'') es donde aparecen los
\emph{orphan nodes}: si un nodo $v \in m(L) \setminus m(K)$ tiene
aristas adyacentes en $G$ que no están en $m(L)$, el pushout
complement no existe y la operación falla. Esto formaliza la noción
clásica de \emph{dangling edges}.

\begin{algorithm}[H]
\SetAlgoLined
\KwIn{$L, K, m, G$}
\KwOut{\textbf{ok} si la gluing condition se cumple, \textbf{fail} si no}
\caption{\textsc{CheckGluing}}
\label{alg:check-gluing}

$\mathrm{ToDelete} \gets m(V_L) \setminus m(V_K)$\;
\ForEach{$v \in \mathrm{ToDelete}$}{
  \ForEach{$e \in E_G$ adyacente a $v$}{
    \If{$e \notin m(E_L)$}{
      
      \Return{\textbf{fail}}\;
    }
  }
}
\Return{\textbf{ok}}\;
\end{algorithm}

\section{Trampas históricas}

\begin{trampa}[Cross-dir rename que crea un ciclo]
\label{tr:rename-cycle}
\citaCommit{(v0.1.x)}. Versión vulnerable: \texttt{rename\_cross\_dir}
no chequeaba que el inode renombrado no fuera ancestro del directorio
destino. \texttt{mv /a /a/b/c} convertía $/a$ en su propio
descendiente, violando \invref{inv:no-dir-cycles}. POSIX manda
\texttt{EINVAL} en ese caso.

\textbf{Fix.} Pre-condición que recorre los \texttt{..} de $d'$ hacia
arriba buscando el inode renombrado; si lo encuentra, retorna
\texttt{InvalidArgument}. Costo $O(\mathrm{depth})$, aceptable.
\textbf{Lección.} Las gluing conditions no son sólo locales (chequear
una arista incidente al matching); pueden requerir traversals globales.
\end{trampa}

\begin{trampa}[Treewidth no acotado por hard links]
\label{tr:link-treewidth}
\citaCommit{(v0.2.0)}. Sin un cap superior a \texttt{link\_count},
\oplink{} permite que un mismo inode acumule miles de aristas
\econtains{} entrantes. Eso degrada las cotas estructurales del
cap.~\ref{cap:topologia} (treewidth lineal en el número de hard
links) y rompe el patrón ``filesystem $\Rightarrow$ tree-like''.

\textbf{Fix.} Constante \texttt{LINK\_MAX = 65535} chequeada en la
gluing condition de \oplink. \textbf{Lección.} Las propiedades
estructurales (treewidth acotado) son tan importantes como las
semánticas; las gluing conditions deben respaldarlas.
\end{trampa}

\section{Resumen del capítulo}

\begin{itemize}
  \item Una regla DPO es un span $\dpo{L}{K}{R}$ con matching $m: L \to
    G$ y gluing condition; el paso construye $H$ vía dos pushouts.
  \item Cada operación POSIX (\opcreate, \opmkdir, \oplink, \opunlink,
    \oprmdir, \oprename, \opwrite) tiene su producción explícita y
    su gluing condition formal.
  \item La gluing condition se verifica antes de mutar; el chequeo de
    invariantes \texttt{check\_invariants} (cap.~\ref{cap:typegraph})
    se verifica después como red de seguridad y como contrato para los
    tests.
  \item En el runtime, las reglas son funciones Rust en
    \citaCrate{sotfs-ops}; en Coq, son teoremas
    \texttt{X\_preserves\_WellFormed} en \citaCoq{Dpo*.v}; en
    TLA\textsuperscript{+}, son acciones del módulo
    \citaTLA{sotfs\_graph.tla}. La correspondencia entre las tres es por
    inspección; cap.~\ref{cap:coq} y \ref{cap:tla} desarrollan cada
    una.
  \item Trampas vistas: orphan blocks bajo unlink hard-linked, ciclo
    cross-dir rename, treewidth sin cap por hard links. Cada una se
    cierra ajustando la gluing condition o agregando una constante
    estructural.
\end{itemize}

\section*{Ejercicios}

\begin{ejercicio}[\dificultad{$\star$}]
\label{ej:dpo-1}
Defina, en una línea, qué es la gluing condition de una regla DPO.
\end{ejercicio}

\begin{ejercicio}[\dificultad{$\star$}]
\label{ej:dpo-2}
Para cada una de las siete operaciones POSIX cubiertas, indique cuál es
el conjunto $K$ (el glue) en términos del subgrafo que sobrevive.
\end{ejercicio}

\begin{ejercicio}[\dificultad{$\star\star$}]
\label{ej:dpo-3}
Diseñe la regla DPO para \texttt{symlink(target, dir, name)} (crear un
enlace simbólico). ¿Qué tipo de nodo introduce? ¿Qué arista? ¿Cuál es
su gluing condition? Compare con \opcreate.
\end{ejercicio}

\begin{ejercicio}[\dificultad{$\star\star$}]
\label{ej:dpo-4}
Lea la regla \texttt{LinkInode} de \citaTLA{sotfs\_graph.tla}. ¿Qué
pasa si el inode $i$ no existe (\texttt{i \textbackslash notin
inodes})? ¿Cómo se compara con el comportamiento POSIX de \oplink{}
sobre un inode borrado?
\end{ejercicio}

\begin{ejercicio}[\dificultad{$\star\star$}]
\label{ej:dpo-5}
Demuestre que \opunlink{} sobre un inode con \texttt{link\_count = 1}
preserva el \invref{inv:no-dangling}. Pista: razonar qué aristas se
borran al desaparecer el inode.
\end{ejercicio}

\begin{ejercicio}[\dificultad{$\star\star$}]
\label{ej:dpo-6}
Formule la NAC de \opmkdir{} como una fórmula de primer orden sobre
$E$. ¿Cuántos chequeos necesita un implementor para verificarla con el
índice \texttt{dir\_name\_idx}?
\end{ejercicio}

\begin{ejercicio}[\dificultad{$\star\star\star$}]
\label{ej:dpo-7}
Diseñe la regla DPO para \texttt{rmdir} sobre un directorio que tenga
\emph{sólo} subdirectorios vacíos (rmdir recursivo). ¿Es una sola regla
DPO? ¿Una secuencia? Justifique en términos de gluing conditions y
atomicidad.
\end{ejercicio}

\begin{ejercicio}[\dificultad{$\star\star\star$}]
\label{ej:dpo-8}
\trampref{tr:rename-cycle} se cierra con un traversal $O(\mathrm{depth})$
sobre los \texttt{..}. Diseñe una versión que use el índice
\texttt{dir\_for\_inode\_idx} para reducir el costo. ¿Sigue siendo
posible un ataque DoS construyendo cadenas profundas?
\end{ejercicio}

\begin{ejercicio}[\dificultad{$\star\star\star$}]
\label{ej:dpo-9}
Cargue el \citaCoq{DpoCreate.v}, lea el theorem
\texttt{create\_preserves\_WellFormed} (línea 335) y reproduzca el
esquema de prueba del lema~\ref{lem:create-preserva} en pseudo-Coq
(intros, casos, cita de los lemmas auxiliares de
\citaCoq{SotfsGraph.v}).
\end{ejercicio}
